% F7 · Dónde vive cada copia de los datos (D12·01)
\documentclass[border=10pt]{standalone}
% ------------------------------------------------------------------
% Preámbulo común de las figuras del curso C1
% Paleta y tipografía alineadas con ../../estilo.tex
% ------------------------------------------------------------------
\usepackage[utf8]{inputenc}
\usepackage[T1]{fontenc}
\usepackage[spanish,es-noquoting]{babel}
\usepackage{helvet}
\renewcommand{\familydefault}{\sfdefault}
\usepackage{xcolor}
\usepackage{tikz}
\usetikzlibrary{arrows.meta,positioning,decorations.pathreplacing,calc,shapes.geometric,fit,backgrounds}

\definecolor{cHead}{RGB}{68,84,106}
\definecolor{cAcc}{RGB}{0,102,153}
\definecolor{cGris}{gray}{0.42}
\definecolor{cSuave}{RGB}{240,244,248}
\definecolor{cBien}{RGB}{34,124,68}
\definecolor{cMal}{RGB}{176,42,42}
\definecolor{cBienSuave}{RGB}{230,243,234}
\definecolor{cMalSuave}{RGB}{250,235,233}
\definecolor{cAccSuave}{RGB}{225,238,246}

\tikzset{
  caja/.style     = {rectangle, rounded corners=2.5pt, draw=cHead, line width=.7pt,
                     fill=cSuave, align=center, inner sep=6pt, font=\small},
  cajaAcc/.style  = {caja, draw=cAcc, fill=cAccSuave},
  cajaOk/.style   = {caja, draw=cBien, fill=cBienSuave},
  cajaNo/.style   = {caja, draw=cMal, fill=cMalSuave, dashed},
  cajaTenue/.style= {caja, draw=cGris, fill=white},
  rotulo/.style   = {font=\scriptsize\itshape, text=cGris, align=center},
  etiq/.style     = {font=\scriptsize, text=cHead, align=center,
                     fill=white, inner sep=1.5pt},
  fl/.style       = {-{Stealth[length=4.5pt]}, draw=cHead, line width=.7pt},
  flAcc/.style    = {fl, draw=cAcc},
  flNo/.style     = {-{Stealth[length=4.5pt]}, draw=cMal, line width=.7pt, dashed},
  llave/.style    = {decorate, decoration={brace, amplitude=4pt, raise=2pt},
                     draw=cGris, line width=.6pt},
  titulo/.style   = {font=\small\bfseries, text=cHead},
}

% Cabecera de figura: #1 número, #2 título
\newcommand{\figtit}[2]{\node[font=\footnotesize\bfseries, text=cGris, anchor=west]}

\begin{document}
\begin{tikzpicture}

\tikzset{empl/.style = {caja, text width=3.8cm, align=left, inner sep=8pt, font=\scriptsize}}

% #1 nombre  #2 quién administra  #3 respaldo
\newcommand{\empl}[3]{%
  {\small\bfseries #1}\\[3pt]
  \textcolor{cGris}{administra:} #2\\
  \textcolor{cGris}{respaldo:} #3}

\node[empl] (cap) at (0,0)
  {\empl{1 · Sistema de captura\\ o de origen}{servicio de informática}{\textcolor{cBien}{diario, 30 días}}};

\node[empl] (tra) at (5.3,0)
  {\empl{2 · Almacenamiento\\ de trabajo}{el grupo}{\textcolor{cMal}{no incluido}}};

\node[empl] (ana) at (10.6,0)
  {\empl{3 · Entorno de análisis}{el grupo}{\textcolor{cMal}{ninguno}}};

\node[empl, draw=cGris, fill=white] (rep) at (16.7,0)
  {\empl{4 · Repositorio}{externo}{política de preservación}};

\draw[fl] (cap) -- (tra);
\draw[fl] (tra) -- (ana);
\draw[fl, draw=cGris, dash pattern=on 2.5pt off 2pt] (ana) -- (rep)
  node[etiq, pos=.5, above=1pt, text=cGris] {al cerrar};

% --- el quinto emplazamiento ----------------------------------------
\node[cajaNo, text width=4.6cm, align=left, inner sep=8pt, font=\scriptsize] (port) at (5.3,-3.3) {
  {\small\bfseries 5 · El portátil de quien analiza}\\[3pt]
  \textcolor{cMal}{no debería existir, y existe casi siempre}\\[2pt]
  sobrevive al proyecto\\ y a veces a la persona};
\draw[flNo] (tra.south) -- (port.north)
  node[etiq, pos=.5, text=cMal] {«lo exporto para el fin de semana»};
\draw[flNo] (ana.south) to[out=-90,in=0] (port.east);

% --- lectura de la figura -------------------------------------------
\draw[llave] ($(cap.north west)+(0,.15)$) -- ($(ana.north east)+(0,.15)$)
  node[midway, above=6pt, font=\small\bfseries, text=cHead]
  {Tres copias vivas a la vez durante el proyecto};

\node[rotulo, text width=16.8cm, align=left, anchor=north west] at (0,-5.1)
  {Cada emplazamiento tiene administrador, respaldo y control de acceso distintos, y el plan los
   declara uno a uno. Lo que queda fuera de la declaración suele ser lo peor protegido: aquí, los
   emplazamientos 2 y 3. Declarar expresamente que el 5 no se usará, y ofrecer la alternativa que
   lo hace innecesario, es lo que convierte una buena intención en una regla del proyecto.};

\end{tikzpicture}
\end{document}
